% Figure: CRT as two sliding rulers --- residues mod 3 and mod 5.
% Source for Figure~\ref{fig:rulers}. TikZ only; regenerable by compiling main.tex.
\begin{tikzpicture}[
  tick/.style={draw=rulegray},
  hit/.style={circle,fill=accent,inner sep=1.6pt},
  lab/.style={font=\footnotesize\sffamily,text=slate},
]
% top ruler: positions 0..14 labeled mod 3
\node[lab,anchor=east] at (-0.3,0) {mod 3:};
\foreach \i in {0,...,14}
{
  \pgfmathtruncatemacro{\v}{mod(\i,3)}
  \draw[tick] (\i*0.55,0) -- (\i*0.55,0.18);
  \node[lab] at (\i*0.55,0.42) {\v};
  \ifnum\v=2 \node[hit] at (\i*0.55,0) {}; \fi
}
% bottom ruler: positions mod 5
\node[lab,anchor=east] at (-0.3,-1.0) {mod 5:};
\foreach \i in {0,...,14}
{
  \pgfmathtruncatemacro{\v}{mod(\i,5)}
  \draw[tick] (\i*0.55,-1.0) -- (\i*0.55,-0.82);
  \node[lab] at (\i*0.55,-0.58) {\v};
  \ifnum\v=3 \node[hit] at (\i*0.55,-1.0) {}; \fi
}
% alignment markers at 8 and 23->position 8
\draw[accent,dashed] (8*0.55,0.15) -- (8*0.55,-1.15);
\node[lab] at (8*0.55,-1.55) {$x = 8$: hits both conditions};
\node[lab] at (8*0.55,-1.95) {(and again at $x = 8 + 15$, $8 + 30$, \dots)};
\end{tikzpicture}
